% SPDX-License-Identifier: Apache-2.0
% covers: Router
\datasheet{Router}{The AXI4 router: one host and N peripherals, decoded by an address map.}

\dsfacts{\code{lib/parts/src/bus/router.rs}}%
{\code{txhdl\_parts::bus::router::Router}}%
{\code{//docs:axi}, \code{//docs:vreteno}, \code{//docs:noc}}

\subsection*{Function}

A router puts several peripherals behind one host. It sits in the five
channels of an AXI4 link, so nothing on either side changes when it is
put there. It takes the host's \code{aw}, \code{ar} and \code{w}, and
gives each peripheral \textit{i} its own, element \textit{i} of the
arrays \code{aws}, \code{ars} and \code{ws}. It takes each peripheral's
answers from \code{bs} and \code{rs}, and gives the host one \code{b}
and one \code{r}. In the netlist element \textit{i} of an array port is
a port of its own, \code{aws\_}\textit{i} and so on.

One unit serves every count of peripherals: the peripheral sides are
arrays of \code{N}, and the lowering unrolls the loops over them when
\code{lowered} runs (issue 500). Each peripheral's range is its entry
of the map, a type \code{M} implementing \code{txhdl::map::AddrMap<N>}
whose constant \code{RANGES} is \code{N} pairs of a base and a mask,
read the same way (issue 593). A burst whose address is in no range is
answered \code{DecErr} by the router itself. The router was a family of
units, \code{Router2} to \code{Router8}, written by \code{router!},
until \#648.

\subsection*{Parameters}

\begin{center}\footnotesize
\begin{tabular}{@{}lp{0.7\textwidth}@{}}
\toprule
Parameter & Meaning \\
\midrule
\code{N} & The number of peripherals. \\
\code{M} & The address map: a type implementing \code{AddrMap<N>},
whose \code{RANGES[i]} is peripheral \textit{i}'s base and mask. It
holds no signal; the router keeps it as \code{PhantomData}. \\
\code{A} & The address width, in bits. \\
\code{D} & The data width, in bits. \\
\code{S} & The write strobe width, one bit per byte lane, so \code{D / 8}. \\
\code{I} & The identifier width, in bits. \\
\bottomrule
\end{tabular}
\end{center}

\subsection*{Address map}

An address \code{a} is peripheral \textit{i}'s when
\code{(a \& mask) == base} for \code{RANGES[i]} and for no range before
it: the first range that matches decides. The router decodes the
address of each address phase. Every address in no range is a hole.
The tables below are generated for the board's router,
\code{BoardRouter} in \code{cpu/vreteno/src/board.rs}, which is
\code{Router<7, BoardMap, 32, 32, 4, 4>} with \code{BoardMap} beside
it:

\begin{center}\footnotesize
\begin{tabular}{@{}llll@{}}
\toprule
Peripheral & Base & Mask & On the board \\
\midrule
0 & \code{0x1000} & \code{0xffff\_f000} & The data memory \\
1 & \code{0x0200\_0000} & \code{0xffff\_0000} & The timer \\
2 & \code{0x3000} & \code{0xffff\_f000} & The AXI-Lite bridge to the
serial port and its neighbours \\
3 & \code{0x4000\_0000} & \code{0xc000\_0000} & The DDR3 memory \\
4 & \code{0x0c00\_0000} & \code{0xfc00\_0000} & The interrupt controller \\
5 & \code{0x0000\_0000} & \code{0xffff\_f000} & The boot memory \\
6 & \code{0x1000\_0000} & \code{0xffff\_0000} & The debug module \\
\bottomrule
\end{tabular}
\end{center}

\dstables{Router}

\subsection*{Behaviour}

\begin{itemize}
\item \textbf{Write address.} The head of \code{aw} is taken when no
write burst is running (\code{wbusy} clear) and the first peripheral
whose range matches has room on its \code{aws} port, or, for a hole,
when no write \code{DecErr} is owed. It is sent to that peripheral
alone.
\item Taking it sets \code{wbusy}, and records the peripheral in
\code{wsel}, a bit each, and the identifier in \code{wid}.
\code{whole} records that no range matched.
\item \textbf{Write beats.} While \code{wbusy} is set, beats pass from
\code{w} to the peripheral in \code{wsel} when it has room. A burst to
a hole has no bit in \code{wsel}, and its beats are taken and dropped.
The beat marked \code{last} clears \code{wbusy}.
\item So a second write address phase, to any peripheral, waits for
the first burst's last beat. AXI4 puts no identifier on \code{w}, and
this keeps two bursts' beats apart.
\item \textbf{Write to a hole.} When its last beat has gone, the router
sets \code{berr} and \code{berr\_id}. It sends a \code{B} of
\code{DecErr} with that identifier when \code{b} has room. That answer
goes before any peripheral's.
\item \textbf{Read address.} The head of \code{ar} is taken when the
first peripheral whose range matches has room on its \code{ars} port,
or, for a hole, when \code{rerr} is clear. Several reads may be
outstanding at once.
\item \textbf{Read to a hole.} Taking it sets \code{rerr}, keeps its
identifier in \code{rerr\_id}, and sets \code{rerr\_left} to its
\code{len} plus one. The router then sends that many beats of
\code{DecErr} and zero data on \code{r}, the last marked \code{last}.
\item \textbf{Read beats.} When no read burst is going up
(\code{rbusy} clear), the router's own error burst goes first. Else the
lowest peripheral that offers a beat is chosen.
\item A peripheral's beat that goes up without \code{last} sets
\code{rbusy} and records the peripheral in \code{rsel}. Its beats then
go up alone until its \code{last}, so \code{last} still marks the end
of one burst.
\item \textbf{Write responses.} When the router owes no \code{DecErr},
the lowest peripheral that offers a \code{B} sends it to the host. Each
response keeps the identifier the peripheral gave it.
\end{itemize}

\subsection*{Verification}

\begin{itemize}
\item Two tests in \code{bus::router} run \code{Router<2, TwoMap, 16,
32, 4, 3>} between a host tracker and two peripheral trackers, in
\code{//lib/parts:parts\_test}. \code{each\_burst\_reaches\_its\_range}
reads from both ranges and reads three beats from a hole, and checks
the hole's answer is \code{DecErr} in three beats.
\code{write\_beats\_do\_not\_interleave} writes three beats to each
peripheral, one burst after the other, and checks each peripheral got
its own burst whole and in order.
\item \code{ex\_axi\_router} issues five bursts before collecting any:
a read to each of two peripherals, a write to the third, and a read and
a write to a hole. The build simulates the netlist of \code{Router<3,
ThreeMap, 16, 32, 4, 3>} against that run under nvc and Verilator:
\code{//docs:axi\_router\_sim\_axi\_router\_tb\_test} and
\code{//docs:axi\_router\_vsim\_test}.
\item Vreteno's run lowers its router, \code{Router<3, RunMap, 32, 32,
4, 2>} in \code{cpu/vreteno/src/main.rs}, and the build simulates that
netlist against the run: \code{//docs:vreteno\_sim\_router\_router\_tb\_test}
and \code{//docs:vreteno\_vsim\_router\_test}.
\code{//cpu/vreteno:lockstep\_test} runs the same router in simulation.
\item \code{four\_peripherals\_share\_one\_node} in \code{bus::noc} puts
a \code{Router<4, NocMap, ..>} behind a node of the network, with a
memory in each range. Two hosts write and read every peripheral and
read a hole.
\end{itemize}

\subsection*{Used by}

\code{Router<3, RunMap, ..>} in Vreteno's run and its lockstep and
pair tests; \code{Router<4, RunMap, ..>}, with the boot memory, in
\code{cpu/vreteno/src/run.rs}; \code{Router<7, BoardMap, ..>} on the
board; \code{Router<4, NocMap, ..>} in the network's test;
\code{Router<3, ThreeMap, ..>} in \code{ex\_axi\_router}; and
\code{Router<2, TwoMap, ..>} in the router's own tests.

\subsection*{Limits}

It serves one host. Several hosts contending for one peripheral are not
handled here, as \code{//docs:axi} states. Write bursts pass one at a
time, whichever peripherals they go to. The merge of responses is by
fixed priority, lowest peripheral first. A mask compares only its own
bits: on an address wider than 16 bits, a mask of \code{0xf000} also
matches its base plus every multiple of \code{0x10000}. Nothing checks
that ranges do not overlap; an address in two ranges goes to the first.
The router decodes on the address alone, and ignores \code{region} and
\code{qos}.
